\begin{textsample}{POS Dim 3 – mt – Score 25.00 – mt/mt\_inf\_22.txt}  \label{ex:f3_pos_084}
2.4.4 As Possibilidades de Experiências com \textbf{Bebês}, \textbf{Crianças} Bem Pequenas e \textbf{Crianças} Pequenas É necessário que o trabalho com \textbf{bebês}, \textbf{crianças} bem \textbf{pequenas} e \textbf{crianças} \textbf{pequenas} proporcione ações didáticas ricas em variadas experiências de comunicação e expressão no cotidiano.
Assim, necessitam de oportunidades de : - Observar o modo como nos relacionamos com os \textbf{bebês} e como eles respondem às nossas ações em todos os momentos : na hora de recebê -los no \textbf{colo} no início do dia ; durante as \textbf{refeições} ; ao longo das atividades ; na hora do banho ; nas \textbf{brincadeiras} ; na troca das \textbf{fraldas} ; - \textbf{Apoiar} as \textbf{crianças} na organização de seus pedidos, apontamentos, observações, etc., favorecendo suas interlocuções no cotidiano ; - \textbf{Conversar} com as \textbf{crianças}, desde \textbf{bebês}, ouvir seus \textbf{balbucios}, decodificar os \textbf{gestos}, expressões faciais, \textbf{sons}, choro, sorriso ; - Ouvir e participar de situações comunicativas no cotidiano ampliando suas referências e aprendendo os usos da linguagem ( expressões faciais, entonação da voz, perguntas e respostas, \textbf{gestos} ) ; - Proporcionar vivências em que as \textbf{crianças} possam expressar -se por meio da linguagem corporal, utilizando movimentos e ações em suas \textbf{brincadeiras}. - Organizar \textbf{pequenos} grupos para \textbf{conversar} com as \textbf{crianças}, apresentar \textbf{brinquedos} e objetos, plantas, animais, fotografias ; - Favorecer às \textbf{crianças} a utilização de recursos midiáticos nos momentos do faz de conta, \textbf{imitação}, fantasia ; - Utilizar, pedagogicamente, diferentes recursos midiáticos ( TV, aparelho telefônico, computador, aparelho de \textbf{som}, dentre outros ) possibilitando que as \textbf{crianças} se expressem oralmente ; - Criar oportunidades para as \textbf{crianças} perguntarem, descreverem e narrarem fatos ; - Estimular a concentração e a participação das \textbf{crianças} em conversas ; - Possibilitar às \textbf{crianças} a expressão de diferentes \textbf{sons} por meio da \textbf{imitação} ; - Organizar e participar com as \textbf{crianças} bem \textbf{pequenas} e \textbf{pequenas} de rodas de conversa e de relatos de episódios pessoais aos \textbf{colegas} da turma ; - Explorar o trabalho com o nome próprio no cotidiano, para marcar pertences, realizar a chamada, criar jogos e \textbf{brincadeiras} ; - Favorecer momentos em que as \textbf{crianças} expressem situações, nomes de objetos e eventos do cotidiano familiar, escolar e de outros contextos de que participam ; - Explorar ao máximo a escrita no cotidiano da instituição ( \textbf{creche} ou escola ), convidando as \textbf{crianças} para participarem tanto da escrita quanto da leitura de bilhetes, comunicados, cardápios, cartazes, agenda do dia, \textbf{rotina}, dentre outras ; - Possibilitar momentos de contação de história, dramatização, \textbf{imitação} de musicalização ; - Organizar momentos, diariamente, de leitura de histórias, desde os \textbf{bebês} ; - Organizar momentos, diariamente, de contação de histórias para as \textbf{crianças}, explorando recursos da oralidade, objetos, fantoches, dedoches, dentre outros ; - Organizar, nas rodas de histórias, momentos em que as \textbf{crianças} possam arriscar -se a \textbf{contar} trechos das histórias aos demais \textbf{colegas}, utilizando sua própria linguagem, com o apoio do professor ; - Organizar momentos de \textbf{brincadeiras} de faz de conta em que as \textbf{crianças} possam arriscar -se a construir narrativas ou partir do repertório de possíveis \textbf{brincadeiras} ( casinha, escritório, castelos e princesas, astronauta, dentre outros ), com a participação do professor ; - Organizar momentos para o uso contextualizado da escrita dos nomes próprios, como, por exemplo, o mural da sala de aula ; - Organizar a biblioteca da sala e assegurar momentos em que as \textbf{crianças} possam explorar os livros ; - Organizar roda de leitura de histórias de diferentes gêneros ; - Organizar rodas de apreciação e indicação de livros e histórias ; - Organizar, nas rodas, momentos em que as \textbf{crianças} possam realizar recontos para os demais \textbf{colegas}, \textbf{apoiando} -se nos livros e colocando em uso os conhecimentos da linguagem escrita ; - Organizar \textbf{pequenos} grupos para produzir textos coletivamente, exercitando assim os conhecimentos da linguagem escrita mesmo antes de saber escrever convencionalmente ; - Organizar momentos de ler e \textbf{brincar} a partir de textos memorizados, tais como letras de música, \textbf{brincadeiras} cantadas, parlendas, poemas, quadrinhas ; - Oferecer materiais para escrever e para ler no contexto das \textbf{brincadeiras} de faz de conta, de acordo com diferentes enredos e temas de \textbf{brincar} ; - Propor situações comunicativas contextualizadas, nas quais escrever se faz necessário para as \textbf{crianças}, \textbf{apoiando} -se no uso de letras móveis ; - Possibilitar às \textbf{crianças} atividades de \textbf{escuta} de história e elaboração de textos de outros gêneros tendo o professor como escriba ; - Selecionar, apresentar, ampliar, pouco a pouco, desde \textbf{bebês}, o repertório de cantigas de ninar e cantigas populares ; - Organizar momentos de cantar e \textbf{brincar} com as \textbf{crianças} desde \textbf{bebês}, ajudando -as a escolher e constituir seu repertório de preferência ; - Organizar o acervo de livros próprios para \textbf{bebês} e os momentos \textbf{diários} para leitura ; - Incluir novos parceiros convidados para as rodas de conversa com os \textbf{bebês}, como, por exemplo, \textbf{crianças} maiores, família, ampliando assim as referências linguísticas dos \textbf{pequenos} ; - Apresentar diferentes produções orais e escritas, variações de \textbf{brincadeiras}, histórias e cantigas, valorizando as diversidades linguísticas regionais e locais ; - Selecionar acervos literários e garantir roda de leitura \textbf{diária} para todas as \textbf{crianças}, a fim de que possam ter uma boa referência, observem e se apropriem dos comportamentos leitores ; - Alimentar os assuntos de interesse das \textbf{crianças} e convidar novos parceiros para as rodas de conversa, ampliando assim o contato das \textbf{crianças} com as variedades linguísticas de sua região. 2.4.5 Observação e Avaliação do Planejamento Assumir o compromisso com as aprendizagens e desenvolvimento das \textbf{crianças} desde a mais tenra idade implica em assegurar determinadas condições didáticas para que a ação pedagógica aconteça intencionalmente, atingindo e alimentando o percurso individual para que ele se desenvolva e amplie.
Assim, observar, \textbf{documentar}, planejar e agir são ações entrelaçadas umas nas outras, de forma dialógica, uma alimentando a outra constantemente, em um círculo virtuoso.
Ao realizar essas ações de forma reflexiva sobre sua prática pedagógica, ou seja, ao tematizar sua prática pedagógica, o professor aprende com a experiência ( como o proposto para as \textbf{crianças} da educação \textbf{infantil} ) e qualifica sua vida profissional.
O planejamento e a avaliação são assim, instrumentos que explicitam a intenção do professor.
Nesse sentido, ao planejar, o professor deverá observar os seguintes aspectos : - Organização de uma \textbf{rotina} \textbf{diária} que observe os critérios de regularidade, continuidade, flexibilidade e diversidade das atividades, vivências e experiências ; - Organização do espaço aconchegante e acolhedor para a roda da leitura, roda da conversa, hora da biblioteca, conto da leitura de memória, de maneira a favorecer interações produtivas entre as \textbf{crianças} ; - Organização do tempo que \textbf{acolha} toda experiência leitora, antes, durante e após a enunciação do texto ; - Previsão e organização na \textbf{rotina}, de momentos para \textbf{conversar} e apreciar o livro/texto após a leitura, estimulando as \textbf{crianças} a escolher os trechos que gostam mais, falando o que sentiram e o que pensaram, estabelecendo relações da história com suas experiências de vida, recomendando ( ou não ) o livro/texto para outros leitores ; - Estabelecimento na \textbf{rotina} \textbf{diária}, do melhor momento para leitura ( antes do parque, depois do almoço ), a partir da experiência das equipes de professores, coordenadores pedagógicos e direção da instituição educativa ; - Detalhamento da leitura : como apresentar o livro selecionado ( por fora e por dentro ), como ajudar a antecipar o conteúdo da história ; como instigar as \textbf{crianças} através de uma pergunta guia ; como preparar as \textbf{crianças} para uma boa \textbf{escuta} da história ; preparar a situação de leitura ( leitura em voz alta, leitura compartilhada, leitura teatralizada ) ; o que \textbf{conversar} com as \textbf{crianças} depois da leitura, o que falar sobre as ilustrações, dentre outras ; - Diversificação do gênero, estilos textuais da literatura \textbf{infantil} para os diversos momentos de leitura presentes na \textbf{rotina} da turma ; - Seleção de livros/textos de qualidade considerando os interesses, as necessidades e os saberes das \textbf{crianças}, colaborando, assim, para a formação de leitores ; - Preparação, pelo professor, da leitura oral e/ou compartilhada apresentando -se as \textbf{crianças} como um bom modelo de leitor e de leitura ; - Previsão e organização de momentos de produção de textos, nos quais o professor é o escriba ; - Previsão e organização de momentos de cantos, músicas, \textbf{escuta} de história, conversas necessárias para o desenvolvimento da oralidade, desde os \textbf{bebês} ; - Organização de materiais que incentivem a interação das \textbf{crianças} com o mundo ficcional ( fantoches, bonecas, livros de literatura ) como incentivo à \textbf{imitação}, sensibilização das \textbf{crianças} ao \textbf{brincar} de ler ( mesmo sem saber ler convencionalmente ), recontar histórias, ao \textbf{brincar} de faz de conta ; - Previsão e organização de tempo para as \textbf{crianças} recontar as histórias que foram lidas para elas ; - Zelo pela expressão oral das narrativas \textbf{infantis}, desde muito cedo, privilegiando o momento \textbf{diário} da jornada para a roda da história ; - Estabelecimento na \textbf{rotina} \textbf{diária}, da roda da conversa.
Algumas orientações “ A cultura escrita, isto é, ações de leitura e de escrita, e objetos portadores de leitura, tais como livros, revistas, jornais, folhetos, histórias em quadrinhos devem estar presentes nas escolas \textbf{infantis} ” ( BARBOSA e OLIVEIRA, 2016, p. 35 ).
O que se pode dizer é que o trabalho com a língua escrita com \textbf{crianças} \textbf{pequenas} não pode decididamente ser uma prática mecânica desprovida de sentido e centrada na decodificação do escrito, conforme recomenda o Parecer n.º 20/2009 das DCNEI ( 2009 ).
Para que o professor possa desempenhar essa importante função de iniciar a formação de leitores e produtores de textos é condição que ele seja leitor e autor de textos.
Nesse sentido, propomos buscar a coerência entre as experiências que devem ser proporcionadas às \textbf{crianças} na Educação \textbf{Infantil} no campo de experiências \textbf{ESCUTA}, FALA, PENSAMENTO E IMAGINAÇÃO, no que se refere às práticas de leitura, e a vivência dessa experiência leitora de literatura pelos professores da Educação \textbf{Infantil}.
Esse processo é denominado de homologia de processos que consiste em experienciar através de todo o processo de formação, as atitudes, modelos didáticos, capacidades e modos de organização que se pretende que venham a ser desempenhados nas práticas pedagógicas com as \textbf{crianças}.
Isso significa que devemos ter com a literatura a mesma relação que propomos que nossas \textbf{crianças} da educação \textbf{Infantil}, desde os \textbf{bebês}, tenham.

% matched lemmas: acolher, apoiar, balbucio, bebê, brincadeira, brincar, brinquedo, colega, colo, contar, conversar, creche, criança, diário, documentar, escuta, fralda, gesto, imitação, infantil, pequeno, refeição, rotina, som
\end{textsample}
